\begin{center}
\textsc{Gauss} an \textsc{Gerling}.\quad Göttingen, 19.~Julius 1827.
\end{center}

\par\noindent\hspace*{\parindent}\mbox{.\ .\ .\ .\ .\ .}\quad Obgleich ich bei meinen astronomischen Operationen nicht in
dem Maasse, wie ich gewünscht hatte, vom Wetter begünstigt bin, so glaube
ich doch die Amplitudo des Bogens zwischen den Sternwarten von Göttingen
und Altona bis auf einen sehr kleinen Bruch einer Secunde festgestellt zu
haben; ich habe eine sehr grosse Menge von Sternen genommen und
zusammen gegen 900 Beobachtungen gemacht. Es bleibt der Unterschied von
\par\clearpage
\noindent\null\hfill{\small\textbf{ZUR HANNOVERSCHEN TRIANGULATION.}}\hfill\makebox[0pt][r]{\fontsize{13}{16}\selectfont 391}\par\medskip
\noindent \(5^{\prime\prime}\) gegen die Rechnung aus den Dreiecken nach \textsc{Walbecks} Erddimensionen,
und es wäre daher sehr interessant gewesen, noch einen Zwischenpunkt zu
haben. Mit dem Sector war es unter den obwaltenden Umständen unmöglich;
vielleicht aber beobachten wir künftig noch mit einem im ersten Vertical
aufzustellenden Passageninstrument gemeinschaftlich in Celle, sowie \textsc{Schumacher}
vorher in Altona und ich nachher in \mbox{Göttingen.\quad\mbox{.\ .\ .\ .\ .\ .}}\quad Noch vor meiner
Reise habe ich auch die sämmtlichen bayerischen Dreiecke mitgetheilt
erhalten, welche durch die Ihrigen mit den meinigen verknüpft \mbox{sind.\quad\mbox{.\ .\ .\ .\ .\ .}}
\par
